\section*{Capítulo \ref{ch:g1:materials-around-us} --- Do que as coisas são feitas}

\begin{solution}{exo:g1:materials-around-us:1}
Uma cadeira é um \omterm{def:g1:materials-around-us:object}{objeto}. A madeira é um \omterm{def:g1:materials-around-us:material}{material}: uma cadeira pode ser feita dela.
\end{solution}

\begin{solution}{exo:g1:materials-around-us:2}
Vidro da janela: vidro. Prego: metal. Pedrinha: pedra. Livro: papel.
Meia: tecido (lã ou algodão).
\end{solution}

\begin{solution}{exo:g1:materials-around-us:3}
Madeira: um lápis, uma porta (ou uma cadeira, uma mesa). Plástico: uma
tigela, um balde (ou um brinquedo, uma garrafa).
\end{solution}

\begin{solution}{exo:g1:materials-around-us:4}
O garfo de metal: é o único que não é feito de madeira nem de plástico.
(Dois são de madeira e um é de plástico.)
\end{solution}

\begin{solution}{exo:g1:materials-around-us:5}
O vidro: ele é transparente.
\end{solution}

\begin{solution}{exo:g1:materials-around-us:6}
O prego de aço. Um ímã puxa o ferro e o aço, não o vidro nem a madeira.
\end{solution}

\begin{solution}{exo:g1:materials-around-us:7}
Dois \omterm{def:g1:materials-around-us:object}{objetos} flutuam (o bloco de madeira e a tampinha de plástico); três
afundam (a pedrinha, a bolinha de gude e o prego).
\end{solution}

\begin{solution}{exo:g1:materials-around-us:8}
As lâminas são de metal, o cabo é de plástico. A tesoura é feita de dois
materiais.
\end{solution}

\begin{solution}{exo:g1:materials-around-us:9}
No aro azul: ela é feita de metal. Ela não está no aro laranja porque
um ímã não a puxa: ela é feita de alumínio, não de ferro nem de aço.
\end{solution}

\begin{solution}{exo:g1:materials-around-us:10}
Uma janela precisa deixar a luz entrar, e precisamos ver através dela.
O vidro é transparente; a madeira não é.
\end{solution}

\begin{solution}{exo:g1:materials-around-us:11}
``Sim'': o prego, o clipe e a chave, 3 \omterm{def:g1:materials-around-us:object}{objetos}. ``Não'': a pedrinha, a
colher, a latinha, a bolinha de gude e a meia, 5 \omterm{def:g1:materials-around-us:object}{objetos}.
$3 + 5 = 8$: cada \omterm{def:g1:materials-around-us:object}{objeto} tem o seu lugar.
\end{solution}
